% !TEX program = xelatex
\documentclass[aspectratio=169]{beamer}
% # TODO:
% French ?
\usepackage{beamerthemeMila}
\usepackage{amsmath}
\usepackage{amssymb}
\usepackage{tikz}
\usepackage{array}
\usepackage[backend=biber,style=numeric,sorting=none]{biblatex}
\addbibresource{references.bib}
\newcommand{\tablebf}[1]{{\addfontfeatures{FakeBold=80}\bfseries #1}}

\title{Mila Theme Template Presentation}
\author{Template Author}
\date{\today}
\milafooter

\begin{document}

\milatitleslide
    {Mila Beamer Theme Template}
    {Date} % Date \today
    {Equation, Vectorized image, and Slide Macro Showcase}

\begin{frame}{Overview}
    \tableofcontents
\end{frame}

\milasectionslide
    {milaTeal}
    {1.0}
    {Mathematics}
    {Complex equations with clear term explanations}

\begin{frame}{Complex Equation Rendering}
\small
\begin{align}
\mathcal{L}(\theta,\phi) =
&\underbrace{\mathbb{E}_{q_{\phi}(z\mid x)}[\log p_{\theta}(x\mid z)]}_{\text{Reconstruction}} \notag - \beta \underbrace{D_{\mathrm{KL}}\!\left(q_{\phi}(z\mid x)\,\|\,p(z)\right)}_{\text{Latent regularization}}
+ \lambda \lVert \theta \rVert_2^2
\end{align}

\begin{notebox}
\textbf{Term legend}
\begin{itemize}
    \item \(x\): observed input data, \(z\): latent variable.
    \item \(q_{\phi}(z\mid x)\): encoder distribution parameterized by \(\phi\).
    \item \(p_{\theta}(x\mid z)\): decoder likelihood parameterized by \(\theta\).
    \item \(D_{\mathrm{KL}}\): divergence to prior \(p(z)\), weighted by \(\beta\).
    \item \(\lambda \lVert \theta \rVert_2^2\): \(L_2\) regularization for parameter stability.
\end{itemize}
\end{notebox}
\end{frame}

\begin{frame}{Results table example}
\footnotesize
\centering
\setlength{\tabcolsep}{4pt}
\begin{tabular}{lcc}
\toprule
\tablebf{Model} & \tablebf{ImageNet (\%) \(\uparrow\)} & \tablebf{MS COCO AP (\%) \(\uparrow\)} \\
\midrule
ResNet-152 & 78.3 & 44.9 \\
ViT-H/14 & \tablebf{90.5} & 58.7 \\
SwinV2-G + Cascade Mask R-CNN & 83.2 & \tablebf{63.1} \\
\bottomrule
\end{tabular}

\vspace{0.6em}
\begin{itemize}
    \item Best values are highlighted in bold on a per-benchmark (column-wise) basis.
\end{itemize}
\end{frame}

\milasectionslide
    {milaOrange}
    {2.0}
    {Visual Content}
    {Raster image rendering and fully vectorized graphics}

\milatwocolumnslide
    {Image Rendering (Raster)}
    {0.55}
    {0.42}
    {
        \includegraphics[width=\linewidth]{\milaPurpleLogoPath}
    }
    {
        \begin{notebox}
            \textbf{Notes}
            \begin{itemize}
                \item Use \texttt{\textbackslash includegraphics} for PNG/JPG assets.
                \item Control placement with relative width (for example \texttt{0.8\textbackslash linewidth}).
                \item Prefer high-resolution source images to avoid blur when projected.
            \end{itemize}
        \end{notebox}
    }

\begin{frame}{Vectorized Rendering (TikZ)}
\centering
\begin{tikzpicture}[>=latex, node distance=3cm]
    \node[draw, rounded corners, fill=milaTealLight, minimum width=2.4cm, minimum height=1cm] (data) {Data};
    \node[draw, rounded corners, fill=milaPurpleLight, right of=data, minimum width=2.8cm, minimum height=1cm] (model) {Model};
    \node[draw, rounded corners, fill=milaGreenLight, right of=model, minimum width=2.6cm, minimum height=1cm] (output) {Predictions};
    \node[draw, rounded corners, fill=milaYellowLight, below of=model, yshift=0.1cm, minimum width=3.6cm, minimum height=1cm] (eval) {Evaluation Metrics};

    \draw[->, thick, milaPurple] (data) -- (model);
    \draw[->, thick, milaPurple] (model) -- (output);
    \draw[->, thick, milaPurple] (output) -- (eval);
    \draw[->, thick, milaPurple] (eval) -- (model);
\end{tikzpicture}

\vspace{0.8em}
\begin{notebox}
Vector graphics (TikZ, pdf, eps, svg, etc.) stay sharp at any zoom level and adapt well to projector scaling.
\end{notebox}
\end{frame}

\milasectionslide
    {milaGreen}
    {3.0}
    {Current Slide Templates}
    {Built-in macros from \texttt{slide\_templates.tex}}

\begin{frame}{Template Catalog}
\begin{notebox}
\textbf{Currently available macros in the theme}
\begin{itemize}
    \item \texttt{\textbackslash milatitleslide\{title\}\{tag\}\{subtitle\}}
    \item \texttt{\textbackslash milasectionslide\{color\}\{label\}\{title\}\{description\}}
    \item \texttt{\textbackslash qasamplelike\{...\}} and \texttt{\textbackslash qasampledislike\{...\}} and \texttt{\textbackslash qasamplewrong\{...\}}
    \item \texttt{\textbackslash milaendslide\{contact\}\{line\}}
    \item \texttt{\textbackslash milatwocolumnslide\{title\}\{left column width\}\{right column width\}\{left column content\}\{right column content\}}
\end{itemize}
\end{notebox}

\vspace{0.4em}
\begin{notebox}
This deck uses all of these templates to showcase the theme.
\end{notebox}
\end{frame}

\qasamplelike
    {Template Example: QA Slide}
    {2026-03-24}
    {Laptop Demo}
    {Can this theme explain a technical result while staying visual?}
    {Yes. Use equation slides for derivations, then pair with image/vector slides for intuition.}
    {This format is useful when presenting model behavior to mixed audiences.}

\qasampledislike
    {Template Example: QA Slide (Dislike)}
    {2026-03-24}
    {Laptop Demo}
    {Can this theme explain a technical result while staying visual?}
    {Yes. Use equation slides for derivations, then pair with image/vector slides for intuition.}
    {This format is useful when presenting model behavior to mixed audiences.}

\qasamplewrong
    {Template Example: QA Slide (Wrong)}
    {2026-03-24}
    {Laptop Demo}
    {Can this theme explain a technical result while staying visual?}
    {Yes. Use equation slides for derivations, then pair with image/vector slides for intuition.}
    {This format is useful when presenting model behavior to mixed audiences.}

\milasectionslide
    {milaPurple}
    {4.0}
    {References}
    {Minimal bibliography example using a .bib file}

\begin{frame}{Selected Papers}
\small
The Transformer architecture and attention-only sequence modeling are introduced in \cite{vaswani2017attention}.

Residual connections that enable very deep neural networks are presented in \cite{he2016resnet}.

Pure Transformer models for large-scale image recognition are detailed in \cite{dosovitskiy2021vit}.
\end{frame}

\begin{frame}{References}
    \printbibliography
\end{frame}

\milaendslide
    {contact@mila.quebec}
    {Thank you for your attention}

\end{document}
